% 本文件由 LaTeX 排版 Agent 根据有效 translation.json 自动生成。
% author=Athanasius; work_id=2818; title=致安提约基雅人书

\chapter{致安提约基雅人书}

% structure: p0001 source=h1 target=omitted reason=page-title-already-rendered-by-parent

致我们亲爱且极渴望的同工\Person{欧瑟伯}、\Person{路西斐}、\Person{阿斯特里乌}、\Person{基马提乌}及\Person{阿纳托利乌}，以及从\Place{意大利}、\Place{阿拉伯}、\Place{埃及}和\Place{利比亚}来到\Place{亚历山大里亚}的主教们：\Person{欧瑟伯}、\Person{阿斯特里乌}、\Person{盖乌}、\Person{阿加图}、\Person{阿摩尼乌}、\Person{阿加托代蒙}、\Person{德拉孔提乌}、\Person{阿德尔菲乌}、\Person{赫尔梅翁}、\Person{马尔谷}、\Person{德奥多鲁}、\Person{安德肋}、\Person{帕弗努提乌}、另一位\Person{马尔谷}、\Person{佐伊鲁}、\Person{梅纳斯}、\Person{乔治}、\Person{路基乌}、\Person{玛加略}及其他人，在基督内向你们致候。\par

我们深信，作为天主的仆役和忠信的管家，你们足以在各方面处理教会的事务。但既然我们听闻，许多先前因嫉妒而与我们分离的人现在渴望和平，同时也有许多人已与阿里乌狂人断绝关系，切愿与我们共融，我们认为最好写信告知你们的善意，就是我们和亲爱的\Person{欧瑟伯}与\Person{阿斯特里乌}所拟定的：你们是我们亲爱且真正最渴望的同工。我们因上述消息而欢欣，并祈祷即使有人仍远离我们，即使有人似乎与阿里乌派意见一致，他都能迅速离开他们的疯狂，以便将来各地所有人都可以说：\BibleQuote{一个主，一个信德} \BibleRef{弗 4:5}。因为正如圣咏作者所说：\BibleQuote{看，弟兄们同居共处，何其善，何其美！} \BibleRef{咏 133:1} 而我们的居所就是教会，我们的心思也应当相同。因为我们相信主也会这样与我们同住，正如他说的：\BibleQuote{我要在他们中间居住，在他们中间行走} \BibleRef{肋 26:12} 和\BibleQuote{这里是我永久的居所，我愿在此居住，因为这是我喜爱的} \BibleRef{咏 132:14}。但\Quote{这里}所指的，不就是那宣讲同一信德和同一宗教的地方吗？\par

2. \Emph{欧瑟伯与阿斯特里乌的使命}\par

我们这些在埃及的人，确实很愿意与我们亲爱的\Person{欧瑟伯}和\Person{阿斯特里乌}一道前往你们那里，原因很多，但主要是为能拥抱你们的友爱，并一同享受所说的和平与和谐。但正如我们在其他信函中所声明的，以及你们可从同工们那里得知的，教会的需要阻碍了我们，我们深表遗憾，便请求我们这些同工\Person{欧瑟伯}和\Person{阿斯特里乌}代替我们到你们那里去。我们感谢他们的虔敬，因为他们虽然本可立即前往自己的教区，却因教会的紧急需要，宁愿不惜一切代价先到你们那里去。因此，他们既已同意，我们便以此自慰：你们和他们在那里，我们全体在心神上与你们同在。\par

3. \Emph{承认\Quote{梅勒提派}，及所有弃绝异端者，尤其关于圣神}\par

因此，凡愿与我们和平的人，特别是那些在旧〔圣堂〕聚会的人，以及那些脱离阿里乌派的人，你们要招来，像父母接纳子女一样接纳他们，像监护人和保护者一样欢迎他们；将他们与我们亲爱的\Person{保利努}及其同伴联合在一起，只要求他们诅咒阿里乌异端，宣认\Place{尼西亚}圣教父们所宣认的信德，并诅咒那些声称圣神是受造物、与基督的\OriginalTerm{本体}{ousia}分离的人。因为拒绝分割\OriginalTerm{天主圣三}{Holy Trinity}，或说圣三的任何部分是受造物，这实际上就是完全弃绝阿里乌派的可憎异端。因为那些人，虽然假装引证\Place{尼西亚}所宣认的信德，却胆敢亵渎圣神，无非是在口头上否定阿里乌异端，内心却依然固守。同时，要让所有人都诅咒\Person{撒伯里乌}和\Person{萨摩撒大的保禄}的邪说，以及\Person{瓦伦提尼安}和\Person{巴西里得}的疯狂，和\Person{摩尼派}的愚妄。因为如果这样做了，所有恶意的猜疑就会从各方面消除，只有天主教会纯洁地彰显出来。\par

4. \Emph{安提约基雅的各方要合一}\par

然而，我们以及那些一直与我们保持共融的人，都持守这信德，我们认为你们中任何人或其他人不会不知。但我们既与所有渴望重新合一的人共同欢欣，特别是与那些在旧〔圣堂〕聚会的人，我们极为光荣主，为一切事，尤其是为这些人的善意，我们劝勉你们在这些条件下与他们建立和谐，并且如我们上面所说的，不附加任何其他条件，也就是说，在旧〔圣堂〕聚会的那些人，或\Person{保利努}及其同伴，不得提出任何其他要求，或任何超出尼西亚定义的事。\par

5. \Emph{撒迪卡信经并非正统信条。\OriginalTerm{自立体}{hypostasis}的问题}\par

并禁止宣读或公布某份文件，即有些人大肆谈论的、据称是在\Place{撒迪卡}会议所拟定的有关信德的文件。因为那次会议并未作出任何此类规定。实际上，当时有些人以尼西亚会议有缺陷为由，要求起草一份信经，并且仓促之中还作了尝试，而在\Place{撒迪卡}聚会的圣公会议对此表示愤慨，并决定不应起草任何信德声明，而应以教父们在\Place{尼西亚}所宣认的信德为满足，因为这份信德毫无缺失，充满虔诚，且不宜颁布第二份信经，以免使\Place{尼西亚}起草的信经被视为不完善，并给那些经常想要起草和定义信经的人提供借口。因此，如果有人提出上述文件或任何其他文件，应加以阻止，并劝他们更要保持和平。因为在这类人身上，我们只看到好争论的动机。至于有些人因某些人谈论三个\OriginalTerm{自立体}{hypostasis}而加以指责，理由是此用语不见于圣经，因而可疑，我们固然认为不应要求超出\Place{尼西亚}宣认之外的任何东西，但由于争议，我们询问了他们，他们是否意指，像阿里乌狂人那样，那些自立体彼此相异、相互外在、在\OriginalTerm{本体}{ousia}上互不相干，且每个自立体各自分离，如同一般受造物，尤其是由人所生的那样，或像不同物质，如金、银、铜那样——或者，是否像其他异端者那样，通过谈论三个自立体，意指三个本原和三个神。\par

他们向我们保证，他们既没有这个意思，也从未持守这种观点。但当我们问他们\Quote{那么你们是什么意思，或者为何用这样的表述？}他们回答说，因为他们信的是一个\OriginalTerm{天主圣三}{Holy Trinity}，不是仅名字上的\OriginalTerm{天主圣三}{trinity}，而是真实存在且自存的，\Quote{我们承认一位真实存在且自存的圣父，一位真实实有且自存的圣子，以及一位自存且真实存在的圣神}，还说他们既没说过有三位神或三个\OriginalTerm{根源}{Beginning}，也绝不容忍任何这么说或这么认为的人，但他们承认一个天主圣三，却只有唯一的\OriginalTerm{天主性}{Godhead}，以及唯一的根源，并且圣子与圣父\OriginalTerm{同性同体}{coessential}，如众教父所说的那样；而圣神既不是受造物，也不是外在的，而是属于且不可分离于圣父和圣子的\OriginalTerm{性体}{Essence}。\par

6. \Emph{关于一个\OriginalTerm{位格}{Hypostasis}或三个位格的问题，不应过分强调。}\par

接受这些人对用语的解释和辩护之后，我们便询问了那些因讲\Quote{一个\OriginalTerm{位格}{Subsistence}}而受他们责备的人，问他们使用这个表述，是像\Person{撒伯里乌}那样，意在否定圣子和圣神，还是仿佛圣子没有本质，或者圣神没有位格。但他们反过来也向我们保证，他们既没有这个意思，也从未持守这种观点，他们说：\Quote{我们用位格这个词，认为说位格和性体是一样的。}又说：\Quote{但我们认为只有一个，因为圣子是出自圣父的性体，并且由于性体的同一性。因为我们相信，只有一个\OriginalTerm{天主性}{Godhead}，且它只有一个性体，而不是圣父有一个性体，圣子和圣神的性体与它有区别。}于是，那些因说有三个位格而受责备的人，便与其他人一致了，而那些讲一个\OriginalTerm{性体}{Essence}的人，也承认了前者的道理（照他们自己的解释）。双方都绝罚了\Person{亚略}，视他为基督的敌人；绝罚了\Person{撒伯里乌}和\Person{撒摩撒达的保禄}，视他们为不敬之人；绝罚了\Person{华伦提努}和\Person{巴西里得}，视他们为背离真理者；绝罚了\Person{摩尼}，视他为制造祸端者。于是在天主的恩宠下，经过上述解释，所有人都一致认为，教父们在\Place{尼西亚}所宣认的信仰比上述这些用语更好，并且今后他们宁愿满足于使用尼西亚的用语。\par

7. \Emph{基督的人性完全，不仅具有肉身。}\par

但是，由于似乎有些人对救主的\OriginalTerm{降生奥迹}{fleshly Economy}争执不休，我们便询问了双方。一方所宣认的，另一方也同意了，那就是：圣言并非像临到先知那样，在时代的末期居住在一个圣人内，而是圣言自己成了血肉，他虽具有天主的形体，却取了奴仆的形体，并为了我们，从\Person{玛利亚}取了肉躯而成为人；这样，人类在他内得以完全且整个地从罪恶中解救出来，从死者中复生，并获赐进入天国。因为他们还宣认，救主并非只有肉身而没有灵魂，也并非没有感觉或理智；因为当主为了我们成为人时，他的肉身不可能没有理智；在圣言自身内完成的救恩也不仅仅是肉身的救恩，也是灵魂的救恩。他既是真实的圣子，也成了人子；既是天主的独生子，同时也成了\BibleQuote{在众多弟兄中作首生者}\BibleRef{罗 8:29}。因此，并非在\Person{亚巴郎}之前有一位天主子，在\Person{亚巴郎}之后有另一位天主子\BibleRef{若 8:58}；也并非有一位复活了\Person{拉匝禄}，另一位却询问关于他的事；而是同一位，作为人说：\BibleQuote{拉匝禄安放在哪里？}作为天主却使他复活了：同一位，作为人且在肉身内吐唾沫，却以天主之子的身份，神性地开了那生来瞎眼的人的眼睛；并且正如\Person{伯多禄}所说\BibleRef{伯前 4:1}，他在肉身内受了苦，却作为天主打开了坟墓，复活了死人。基于这些理由，他们这样理解福音中所说的一切，并向我们保证，对于圣言的\OriginalTerm{降生成人}{Incarnation}和成为人，他们持守同样的真理。\par

8. \Emph{不可让用词问题分裂想法相同的人。}\par

这些既已如此宣认，我们劝你们不要草率地谴责那些这样宣认并如此解释自己用语的人，也不要拒绝他们，反而要接纳他们，因为他们渴望和平且为自己辩护；但同时，对于那些拒绝这样宣认并解释其用语的人，你们要查禁并斥责，视他们为持有可疑观点。你们虽然不宽容后者，但也要劝告那些解释并持守正确观点的人，不要进一步探究彼此的意见，也不要为词句而徒然争辩，不要再为上述用语继续争执，而应在虔诚的心意上达成一致。因为那些不这么想，只是用这类琐碎用语挑起纷争，并寻求超越\Place{尼西亚}所拟定内容的人，所做的无非是\BibleQuote{把混浊的混乱给近人喝}\BibleRef{哈 2:15}，就像那些怨恨和平、喜爱分裂的人。但你们，作为善良的人、忠信的仆人和主的管家，要制止并杜绝那些引起反感且怪异的事，在信仰纯正的前提下，视此类和平高于一切。或许天主会怜悯我们，使分裂的联合起来，并使羊群再度合成一个\BibleQuote{一个羊群}\BibleRef{若 10:16}，我们众人同属一个领袖，即我们的\Person{主耶稣基督}。\par

9. \Emph{上述用语已获一致同意。}\par

对于这些事，虽然原无需要求任何超过尼西亚大公会议的内容，也不应容忍争辩的言辞，但为了和平，并避免那些愿意正确信仰的人被拒绝，我们仍进行了查问。他们承认的，我们简要地记录下来，由我们这些留在\Place{亚历山大城}的人，会同我们的同工\Person{阿斯特里乌斯}和\Person{欧瑟比乌斯}执笔。因为我们大多数人已经各自回自己的教区去了。你们要在你们惯常聚会的地方公开宣读这封信，并请乐意邀请所有人前来。因为这封信理当先在那里宣读，然后那些渴望并致力于和平的人应当在那里重新合一。再后，当他们重新合一之后，在众平信徒认为最适宜的地点，当着你们各位的面，举行公众集会，同声光荣主。与我在一起的弟兄们问候你们。我祈求你们安康，并在主前记念我们；我，\Person{亚大纳削}，以及一同聚集的主教们，也都签署；还有\Place{撒丁岛}主教\Person{路基斐尔}所派遣的两位执事，\Person{赫伦尼乌斯}和\Person{阿伽佩图斯}；\Person{保利诺}那边的\Person{马克西莫斯}和\Person{卡勒梅鲁斯}两位执事。此外还有主教\Person{阿波利纳里乌斯}的一些隐修士，奉他差遣为此事而来。\par

10. \Emph{署名}\par

这封信所寄给的主教们的名字如下：\Place{高卢}的\Place{维尔吉利}城的\Person{欧瑟比乌斯}、\Place{撒丁岛}的\Person{路基斐尔}、\Place{阿拉伯}\Place{培特拉}的\Person{阿斯特里乌斯}、\Place{科厄勒叙利亚}\Place{帕尔图斯}的\Person{基玛提乌斯}、\Place{欧波亚}的\Person{安纳托利乌斯}。\par

发信人：教宗\Person{亚大纳削}，及在\Place{亚历山大城}与他一起的人，即：\Person{欧瑟比乌斯}、\Person{阿斯特里乌斯}和上述其他人，还有\Place{利比亚近区}的\Place{帕拉托尼乌姆}的\Person{该尤斯}、\Place{埃及}\Place{埃勒阿尔基亚}和\Place{弗拉戈尼斯}地区的\Person{阿加图斯}、\Place{帕克内穆尼斯}和\Place{埃勒阿尔基亚}其余地区的\Person{阿摩尼乌斯}、\Place{斯凯迪亚}和\Place{梅内拉伊塔斯}的\Person{阿加托德蒙}、\Place{小赫尔穆波利斯}的\Person{德拉孔提乌斯}、\Place{利克尼}的\Place{奥努菲斯}的\Person{阿德尔菲乌斯}、\Place{塔内斯}的\Person{赫尔弥翁}、\Place{利比亚近区}\Place{齐格拉}的\Person{玛尔库斯}、\Place{阿特里比斯}的\Person{狄奥多鲁斯}、\Place{阿尔塞诺埃}的\Person{安德烈亚斯}、\Place{萨伊斯}的\Person{帕弗努提乌斯}、\Place{菲莱}的\Person{玛尔库斯}、\Place{安德罗斯}的\Person{佐伊卢斯}、\Place{安提弗拉}的\Person{梅纳斯}。\par

\Person{欧瑟比乌斯}也用拉丁文签署了以下内容，其译文如下：\par

\Signature{我\Person{欧瑟比乌斯}，根据你们双方就\OriginalTerm{位格}{Subsistences}达成一致的确切声明，也加入我的同意；此外，关于我们救主的\OriginalTerm{降生成人}{Incarnation}，即\Person{天主子}成为人，完全取了人性而无罪过，如同我们旧人的构成，我批准信函的文本。再者，由于排除了\WorkTitle{萨尔底卡文书}，以免有超出\WorkTitle{尼西亚信经}之嫌，我也加入我同意，为使\WorkTitle{尼西亚信经}不因此被视为遭到排拒，并同意不应发布该文书。我在主内祈求你们安好。}\par

\Signature{我\Person{阿斯特里乌斯}赞同以上所写，并在主内祈求你们安好。}\par

11. \Emph{在\Place{安提约基雅}签署的\WorkTitle{通谕}}。\par

在这\WorkTitle{通谕}由\Place{亚历山大城}发出，并经上述人士签署之后，收信人也依次签署如下：\par

\Signature{我\Person{保利诺}持守如此信仰，正如我从教父们所领受的：\Person{圣父}完美地存在并自立，\Person{圣子}完美地自立，\OriginalTerm{圣神}{Holy Spirit}也完美地自立。因此，我也接受上文关于三个\OriginalTerm{位格}{Three Subsistences}，及一个\OriginalTerm{位格}{one Subsistence}、或更好说一个\OriginalTerm{本质}{Essence}的解释，并接受如此持守的人。因为虔敬地持守并宣认在同一\OriginalTerm{神性}{Godhead}内的\OriginalTerm{天主圣三}{Holy Trinity}是合宜的。关于\Person{圣父}的\OriginalTerm{圣言}{Word}为我们成为人，我持守经上所记，即如\Person{若望}所说：\BibleQuote{圣言成了血肉}，并非如同那些极不虔敬者所说的，谓他经历了变化，而是他为了我们成为人，由\OriginalTerm{圣神}{Holy Spirit}和\Person{童贞玛利亚}所生。因为我们的救主拥有的身体，并非没有灵魂、没有感觉、没有理智。因为\OriginalTerm{主}{Lord}为我们成为人，若他的身体没有理智，这是不可能的。因此，我绝罚那些弃绝\WorkTitle{尼西亚信经}所宣认的信仰，以及不说\Person{圣子}是由\Person{圣父}的\OriginalTerm{本质}{Essence}而来、与\Person{圣父}\OriginalTerm{同一性体}{coessential}的人。再者，我绝罚那些说\OriginalTerm{圣神}{Holy Spirit}是经由\Person{圣子}所造成的\OriginalTerm{受造物}{Creature}的人。我再次绝罚\Person{萨培理}和\Person{福蒂诺}的异端，以及一切异端，行在\WorkTitle{尼西亚信经}信仰及上文所写的一切中。}\Signature{我\Person{卡尔特里乌斯}祈求你们安好。}\par

\SourceRecord{Translated by H. Ellershaw.；Nicene and Post-Nicene Fathers, Second Series；Vol. 4.；Edited by Philip Schaff and Henry Wace.；Buffalo, NY: Christian Literature Publishing Co.,；1892.；<http://www.newadvent.org/fathers/2818.htm>.}
